\documentclass[10pt,a4paper]{amsart}
\usepackage[margin=19mm]{geometry}
\usepackage{amsmath,amssymb,mathtools}
\usepackage[scr=rsfs,cal=euler]{mathalfa}
\usepackage{xfrac}
\usepackage[unicode,hidelinks,bookmarks=false]{hyperref}
\newcommand{\ie}{i.e.\ }
\newcommand{\cf}{cf.\ }
\newcommand{\resp}{respectively\ }
\newcommand{\Subsection}{\subsection}
\providecommand{\nobreakdash}{\nobreak\hbox{-}}
\setlength{\emergencystretch}{2em}
\multlinegap0pt
\usepackage{mathtools}
\usepackage{amssymb}
\usepackage[shortlabels]{enumitem}
\usepackage{tikz}
\newtheorem{theorem}{Theorem}
\newcommand{\dd}{\mathrm{d}}
\newcommand{\ee}{\mathrm{e}}
\newcommand{\ii}{\mathrm{i}}
\title{The Benjamin-Ono Equation in the Long-Time Limit: Linearized Self-Similar Universality}
\author{Louise Gassot, Patrick G\'erard,, Peter D. Miller}
\date{Source: arXiv:2606.02989v1 (CC BY 4.0)}
\begin{document}
\maketitle

\noindent\emph{Source and licence.} The Benjamin-Ono Equation in the Long-Time Limit: Linearized Self-Similar Universality. Version arXiv:2606.02989v1 (CC BY 4.0), distributed by arXiv under the Creative Commons Attribution 4.0 International licence, http://creativecommons.org/licenses/by/4.0/, as recorded in the arXiv licence metadata for this exact version and shown on the abstract page https://arxiv.org/abs/2606.02989v1. Full licence notice: \texttt{licenses/arxiv-2606.02989-CC-BY-4.0.txt}.\par\smallskip

\noindent\textit{Editorial scope.} Counted unit: the article's theorem theorem (source0.tex lines 331--333). The article's complete original proof of it is delivered with it (source0.tex lines 334--377, one \texttt{proof} environment of 44 lines). No mathematics is written, repaired or re-derived: the statement and the proof are the article's own lines, in the article's own notation, and no retained line is substituted or re-flowed.  Retained uncounted as the source-local context the proof needs: lines 275--278; lines 280--283.  Nothing was omitted from any retained span, so the package carries no omission record.  No retained line contains a citation, so the wrapper carries no bibliography and no bibliography entry.  No retained line contains a figure environment, an image-inclusion command or drawing code, so no image is delivered and the package declares no asset dependency.  Editorial formatting only, in addition: an AMS \texttt{amsart} preamble that restates the article's own declarations and macros used by the retained lines, namely the environments \texttt{theorem} on separate counters, together with the standard packages those lines need.  The retained environments print in this document as Theorem 1 in order of appearance, whereas the article prints the same items with the numbering of its own full document.  The complete native source of the article as distributed in the arXiv e-print for this exact version is bundled unchanged as \texttt{sources/201947/source0.tex}.
\begin{equation}
\mathrm{Re}(c_1+\cdots+c_N)=0.      
\label{eq:XS-u0-in-L1}
\end{equation}

\begin{equation}
        u_0(x)=\sum_{n=1}^N\left[\frac{c_n}{x-p_n} + \frac{c_n^*}{x-p_n^*}\right],\quad x\in\mathbb{R}.
    \label{eq:XS-u0-rational}
\end{equation}

\begin{theorem}
    Potentials $u_0$ of the form \eqref{eq:XS-u0-rational} satisfying \eqref{eq:XS-u0-in-L1} and for which $\Delta\neq 0$ are dense in $L^2(\mathbb{R},\mathbb{R})$.
\end{theorem}

\begin{proof}
Let $v\in L^2(\mathbb{R},\mathbb{R})$ and $\epsilon>0$ be given.  We will find a potential $u_0$ with the desired properties so that $\|u_0-v\|_{L^2(\mathbb{R})}<\epsilon$.  

    First note that simple-pole rational functions $u_0$ of the form \eqref{eq:XS-u0-rational} satisfying the condition \eqref{eq:XS-u0-in-L1} and having no exceptional indices are dense in $L^2(\mathbb{R},\mathbb{R})$.  Therefore, there exists such a function $u_0^\epsilon$ for which $\|u^\epsilon_0-v\|_{L^2(\mathbb{R})}<\frac{1}{2}\epsilon$.  If $\Delta\neq 0$ for $u_0^\epsilon$, then we are done because $\|u_0^\epsilon-v\|_{L^2(\mathbb{R})}<\frac{1}{2}\epsilon<\epsilon$.
    
    Otherwise, we will find another simple-pole rational function $\widetilde{u}_0^\epsilon$ of the same type but for which $\Delta\neq 0$ and for which $\|\widetilde{u}_0^\epsilon-u_0^\epsilon\|_{L^2(\mathbb{R})}<\frac{1}{2}\epsilon$ so that $\|\widetilde{u}_0^\epsilon-v\|_{L^2(\mathbb{R})}\le \|\widetilde{u}_0^\epsilon-u_0^\epsilon\|_{L^2(\mathbb{R})}+\|u_0^\epsilon-v\|_{L^2(\mathbb{R})}<\frac{1}{2}\epsilon+\frac{1}{2}\epsilon=\epsilon$.  Fix $N=N(\epsilon)$ as the number of poles in the upper half-plane for $u_0^\epsilon$, and fix also the poles $p_n=p_n^\epsilon$, $n=1,\dots,N$.  We will construct $\widetilde{u}_0^\epsilon$ by perturbing the coefficients $c_n=c_n^\epsilon$ of $u_0^\epsilon$ consistent with the condition \eqref{eq:XS-u0-in-L1}.  Let $\mathbb{H}$ denote the complex hyperplane in $\mathbb{C}^{2N}$ consisting of vectors $\mathbf{v}$ with components satisfying $v_1+v_3+\cdots + v_{2N-1}=0$. The identification $c_n=v_{2n-1}+\ii v_{2n}$ and $c_n^*=v_{2n-1}-\ii v_{2n}$ shows that if $\mathbf{v}$ is real, then $\mathbf{v}\in\mathbb{H}$ is the same as the condition \eqref{eq:XS-u0-in-L1}.  Now allowing $\mathbf{v}\in\mathbb{H}$ to be complex,  we define an entire function $\widetilde{\Delta}:\mathbb{H}\to\mathbb{C}$ by applying the same procedure used to define $\Delta$ except that we define the contours $C_1,\dots,C_N$ and the quantities $E_1,\dots,E_N$ as if all indices $n=1,\dots,N$ were non-exceptional.  Clearly, if $\mathbf{v}$ is real, then $\widetilde{\Delta}=\Delta$ for fixed $p_1,\dots,p_N$, unless there is an exceptional index.  Since no indices are exceptional for $u_0^\epsilon$, which is an open condition on the coefficients $c_1,\dots,c_N$, there is a neighborhood $U$ of the real part $\mathbb{H}_\mathbb{R}$ of $\mathbb{H}$ containing $\mathbf{v}_0:=(\mathrm{Re}(c_1^\epsilon),\mathrm{Im}(c_1^\epsilon),\dots,\mathrm{Re}(c_N^\epsilon),\mathrm{Im}(c_N^\epsilon))^\top$ on which $\Delta=\widetilde{\Delta}$ is real-analytic.
    
    We claim that $\widetilde{\Delta}$ does not vanish identically on $\mathbb{H}$.  To see this, we choose a point of the form $\mathbf{v}=(0,v_2,0,\dots,0)^\top\in\mathbb{H}$ (or $c_2=\cdots=c_N=0$ while $c_1=\ii v_2$ and $c_1^*=-\ii v_2$), so that $\ee^{-\ii L(z)}$ is the branch of $\ee^{-\ii L(z)}=[(z-p_1)/(z-p_1^*)]^{v_2}$ that tends to $1$ as $z\to\infty$ backward along each contour $C_1,\dots,C_N$. Then, $E_1=E_2=\cdots=E_N=\ee^{2\pi\ii v_2}$ is the (common) corresponding limiting value along each contour in the forward direction, so all elements of the vector $\mathbf{d}$ are the same:  $d_m=1-E_m=1-\ee^{2\pi\ii v_2}$, $m=1,\dots,N$.  In the matrix elements of $\mathbf{J}$, the initial contour arcs $C_m^-$ may all be taken to be the same (denoted $C^-$), while a residue computation shows that
    \begin{equation}
        \int_{C_m^+}\frac{\ee^{-\ii L(z)}-\ee^{2\pi\ii v_2}}{z-p_n}\,\dd z = \int_{C_1^+}\frac{\ee^{-\ii L(z)}-\ee^{2\pi\ii v_2}}{z-p_n}\,\dd z + 2\pi\ii (\ee^{-\ii L(p_n)}-\ee^{2\pi\ii v_2})\delta_{1<n\le m},\quad m\ge 2.
    \end{equation}
    Therefore, setting 
    \begin{equation}
        f_n:=\int_{C^-}\frac{\ee^{-\ii L(z)}-1}{z-p_n}\,\dd z +\int_{C_1^+}\frac{\ee^{-\ii L(z)}-\ee^{2\pi\ii v_2}}{z-p_n}\,\dd z,\quad n=1,\dots,N,
    \end{equation}
    we see that the first row elements of $\mathbf{J}$ are
    $J_{1n}=f_n+2\pi\ii\ee^{2\pi\ii v_2}\delta_{n=1}$ while the subsequent rows have elements that can be written in the form $J_{mn}=J_{1n}+2\pi\ii\ee^{-\ii L(p_n)}\delta_{1<n\le m}$ for $m=2,\dots,N$.  Therefore, by row operations subtracting the first row from each subsequent row,
    \begin{multline}
        \det(\mathbf{J}\mathop{\longleftarrow}^n\mathbf{d}) \\= 
        \det\left(\begin{bmatrix}f_1+2\pi\ii\ee^{2\pi\ii v_2} & f_2 & f_3 & 
        \cdots & f_N\\        
        0 & 2\pi\ii\ee^{-\ii L(p_2)} & 0 & 
        \cdots & 0\\        
        0 & 2\pi\ii\ee^{-\ii L(p_2)} & 2\pi\ii\ee^{-\ii L(p_3)} & 
        \ddots & \vdots\\           
        \vdots & \vdots & \vdots &
        \ddots & 0\\        
        0 & 2\pi\ii\ee^{-\ii L(p_2)}& 2\pi\ii\ee^{-\ii L(p_3)} & 
        \cdots &2\pi\ii \ee^{-\ii L(p_N)}\end{bmatrix}\mathop{\longleftarrow}^n\begin{bmatrix}1-\ee^{2\pi\ii v_2} \\0\\0\\
        \vdots\\ 0\end{bmatrix}\right).
    \end{multline}
    This determinant obviously vanishes unless $n=1$, in which case it is easily evaluated.  We deduce that
    \begin{equation}
        \widetilde{\Delta}(0,v_2,0,\dots,0)=
    \det(\mathbf{J}\displaystyle\mathop{\longleftarrow}^1\mathbf{d}) = (1-\ee^{2\pi\ii v_2})(2\pi\ii)^{N-1}\prod_{n=2}^N\ee^{-\ii L(p_n)},
    \end{equation}
    which only vanishes if $v_2\in\mathbb{Z}$.  

    Returning to the neighborhood $U\subset\mathbb{H}_\mathbb{R}$, the real and imaginary parts of $\Delta=\widetilde{\Delta}$ are both real-analytic functions on $U$.  Therefore $m:=|\Delta|^2=\mathrm{Re}(\Delta)^2+\mathrm{Im}(\Delta)^2$ is a real-analytic function on $U$ that does not vanish identically, although $m=0$ for $\mathbf{v}=\mathbf{v}_0\in U$.  We introduce explicit coordinates $\mathbf{x}=(x_1,\dots,x_{2N-1})^\top$ on $U$ with $\mathbf{v}=\mathbf{v}_0$ corresponding to $\mathbf{x}=\mathbf{0}$ by 
    \begin{equation}
       \mathbf{v}=\mathbf{v}(\mathbf{x}):=\mathbf{v}_0+\left(-\sum_{j=1}^{N-1}x_{2j},x_1,x_2,\dots,x_{2N-1}\right)^\top 
    \end{equation}
    and we write $m=m(\mathbf{x})$.  Let $k\ge 1$ be the smallest integer such that $\partial_\mathbf{x}^\alpha m(\mathbf{0})=0$ for all multi-indices $\alpha$ with $|\alpha|<k$.  Note that $k<\infty$ because $m(\mathbf{x})$ does not vanish identically.  Let $P_k(\mathbf{x})$ denote the (nonvanishing) homogeneous polynomial consisting of the terms of degree $k$ in the Taylor expansion of $m(\mathbf{x})$ about $\mathbf{x}=\mathbf{0}$.  Then there is a unit vector $\mathbf{e}\in\mathbb{R}^{2N-1}$ such that $P_k(\mathbf{e})\neq 0$.  For $\rho>0$, we then have $m(\rho\mathbf{e})=P_k(\rho\mathbf{e})+O(\rho^{k+1})=\rho^kP_k(\mathbf{e}) + O(\rho^{k+1})$ as $\rho\downarrow 0$.  In particular, $m(\rho\mathbf{e})\neq 0$ for $\rho>0$ sufficiently small.  Since for fixed $N$ and $p_1,\dots,p_N$, the rational form \eqref{eq:XS-u0-rational} defines a continuous map from the coefficients $c_1,\dots,c_N$ into $L^2(\mathbb{R})$, choosing $\widetilde{\mathbf{v}}=\mathbf{v}(\mathbf{x})=\mathbf{v}(\rho\mathbf{e})$ for $\rho>0$ sufficiently small gives the coefficients $\widetilde{c}_n=\widetilde{v}_{2n-1}+\ii \widetilde{v}_{2n}$ of a rational function $\widetilde{u}_0^\epsilon$ with the desired properties, and the proof is complete.
\end{proof}

\end{document}
